\section{Hastings' approximate factorization of the ground-state projector}
\label{sec:hastings_lemma1}

% Hastings, arXiv:0705.2024v4, law16.tex: statement of Lemma 1 at lines 246--267,
% proof in the appendix at lines 763--910, note added on positivity at lines 912--1003.

Lemma~1 of \cite{Hastings2007AreaLaw} approximates the ground-state projector
$P_0$ of a gapped finite-range chain by a product $O_BO_LO_R$ of contractions
supported on the $2l$ sites around a cut $j$, the left sites $X_{1,j}$, and
the right sites $X_{j+1,N}$. Following the appendix of
\cite{Hastings2007AreaLaw} and the three-way Hamiltonian splitting of
Section~\ref{sec:hastings_three_way_splitting}, let $H=H^\dagger$ have ground
energy $E_0$ and spectral gap $\Delta E>0$, with operator norm $\|\cdot\|$.
Under the index shift $\{1,\dots,N\}\mapsto\{0,\dots,L-1\}$ ($N=L$) of
Section~\ref{sec:hastings_three_way_splitting}, $X_{1,j}=\{0,\dots,j-1\}$ and
$X_{j+1,N}=\{j,\dots,L-1\}$. For the Lieb--Robinson parameters $(v,\xi_C)$ of
\cite[lines 96--112]{Hastings2007AreaLaw}, \cite[lines 124--144]{Hastings2007AreaLaw} sets
\begin{align}
    \xi & =\max(2v/\Delta E,\xi_C)
    \label{eq:hastings_lemma1_xi}
\end{align}
and $\xi'=6\xi$. For fixed $\xi$, the appendix argument yields every rate
$\xi'>6\xi$ rather than $\xi'=6\xi$; below $\xi'>6\xi$, except in
Lemma~\ref{lem:hastings_lemma1} ($\xi'=6\xi$ as printed) and
Theorem~\ref{thm:hastings_lemma1_endpoint} ($\xi'=6\xi$ for a raised velocity).

Henceforth $H_L,H_B,H_R$ denote the centered parts
$H_\bullet-\langle\Psi_0,H_\bullet\Psi_0\rangle$ of
Section~\ref{sec:hastings_three_way_splitting}
\cite[lines 781--782]{Hastings2007AreaLaw}. Since scalar multiples of $1$
commute with $H$ and are fixed by the Gaussian filter, centering leaves
$[H,H_\bullet]$ and $\widetilde H_\bullet-H_\bullet$ unchanged.
% Lean: TNLean.Lattice.gaussianFilteredTerm_sub_smul_one and
% TNLean.Lattice.gaussianFilteredTerm_sub_smul_one_sub (filtering fixes scalars).

\subsection*{Gaussian-filtered operators}

\begin{definition}[Gaussian-filtered operator]
    \label{def:hastings_gaussian_filtered_term}
    \lean{TNLean.Lattice.gaussianFilteredTerm}
    \leanok
    For $q>0$ and a matrix $W$, the \emph{Gaussian-filtered operator} is
    \begin{align}
        \widetilde W & =\frac{\Delta E}{\sqrt{2\pi q}}\int_{-\infty}^{\infty}
        e^{iHt}We^{-iHt}\,e^{-(t\Delta E)^2/(2q)}\,dt,
        \label{eq:hastings_gaussian_filtered_term}
    \end{align}
    and $\widetilde W=W$ for $q=0$. Applied to the centered parts
    $W\in\{H_L,H_B,H_R\}$, (\ref{eq:hastings_gaussian_filtered_term}) gives
    $\tilde H_L^0,\tilde H_B^0,\tilde H_R^0$ of
    \cite[lines 783--793]{Hastings2007AreaLaw}.
\end{definition}

\begin{theorem}[Spectral properties of filtered operators]
    \label{thm:hastings_gaussian_filtered_spectral}
    \lean{
        TNLean.Lattice.inner_gaussianFilteredTerm,
        TNLean.Lattice.inner_gaussianFilteredTerm_ground,
        TNLean.Lattice.norm_gaussianFilteredTerm_le,
        TNLean.Lattice.gaussianFilteredTerm_sum,
        TNLean.Lattice.gaussianFilteredTerm_self,
        TNLean.Lattice.sum_gaussianFilteredTerm,
        TNLean.Lattice.mem_groundEigenspace_iff,
        TNLean.Lattice.inner_eq_zero_of_mem_groundEigenspace,
        TNLean.Lattice.norm_sub_starProjection_gaussianFilteredTerm_le
    }
    \leanok
    \uses{def:hastings_gaussian_filtered_term}
    If $Hv=Ev$ and $Hw=E'w$, then
    \begin{align}
        \langle v,\widetilde Ww\rangle
         & =e^{-q((E-E')/\Delta E)^2/2}\,\langle v,Ww\rangle,
        \label{eq:hastings_gaussian_filtered_matrix_element}
    \end{align}
    so $\langle v,\widetilde Ww\rangle=\langle v,Ww\rangle$ when $E=E'=E_0$.
    The map $W\mapsto\widetilde W$ is linear, satisfies
    $\|\widetilde W\|\le\|W\|$ and $\widetilde H=H$, and thus
    $\sum_i\widetilde H_i=H$ whenever $\sum_iH_i=H$
    \cite[line 793]{Hastings2007AreaLaw}. For a ground vector $\psi$ and the
    ground-space projection $P_{\mathrm{gs}}$,
    $\|(1-P_{\mathrm{gs}})\widetilde W\psi\|\le\|W\|e^{-q/2}\|\psi\|$.
\end{theorem}
\begin{proof}\leanok
    In the $H$-eigenbasis, $\langle v,e^{iHt}We^{-iHt}w\rangle
        =e^{i(E-E')t}\langle v,Ww\rangle$, whose Gaussian integral is
    (\ref{eq:hastings_gaussian_filtered_matrix_element}). Unitaries preserve
    $\|W\|$, and $e^{iHt}He^{-iHt}=H$. Since $|E-E_0|\ge\Delta E$ on the
    excited subspace, each excited coefficient of $\widetilde W\psi$ is scaled
    by at most $e^{-q/2}$.
\end{proof}

\begin{theorem}[Filtered operators without ground components]
    \label{thm:hastings_filtered_no_ground_component}
    \lean{
        TNLean.Lattice.norm_gaussianFilteredTerm_apply_le_commutator,
        TNLean.Lattice.eigenvalues_eq_or_le_of_isGappedGroundState,
        TNLean.Lattice.norm_gaussianFilteredTerm_apply_le_commutator_of_isGappedGroundState
    }
    \leanok
    \uses{thm:hastings_gaussian_filtered_spectral}
    If $H\psi=E_0\psi$ and $P_{\mathrm{gs}}W\psi=0$, then
    \begin{align}
        \|\widetilde W\psi\| & \le\Delta E^{-1}e^{-q/2}\,\|[H,W]\psi\|.
        \label{eq:hastings_filtered_no_ground_component}
    \end{align}
    In particular, (\ref{eq:hastings_filtered_no_ground_component}) holds for a
    unique gapped ground state $\Psi_0$ when $\langle\Psi_0,W\Psi_0\rangle=0$,
    hence for the centered parts $H_L,H_B,H_R$
    \cite[lines 794--815]{Hastings2007AreaLaw}.
\end{theorem}
\begin{proof}\leanok
    For an $H$-eigenvector $v_E$ with eigenvalue $E$,
    $\langle v_E,[H,W]\psi\rangle=(E-E_0)\langle v_E,W\psi\rangle$. Since
    $\langle v_{E_0},W\psi\rangle=0$ and $E-E_0\ge\Delta E$ on the excited
    subspace, (\ref{eq:hastings_gaussian_filtered_matrix_element}) gives
    $|\langle v_E,\widetilde W\psi\rangle|
        \le\Delta E^{-1}e^{-q/2}|\langle v_E,[H,W]\psi\rangle|$.
\end{proof}

\begin{theorem}[Locality of filtered operators]
    \label{thm:hastings_filtered_locality}
    \lean{
        TNLean.Lattice.norm_map_gaussianFilteredTerm_le_mass,
        TNLean.Lattice.commutatorCLM,
        TNLean.Lattice.norm_commutator_gaussianFilteredTerm_le_mass,
        TNLean.Lattice.gaussianReal_real_compl_Icc_le_exp,
        TNLean.Lattice.gaussianFilter_tail_le,
        TNLean.Lattice.norm_commutator_gaussianFilteredTerm_le
    }
    \leanok
    \uses{def:hastings_gaussian_filtered_term}
    If $q>0$, $T,\varepsilon\ge0$, and
    $\|[e^{iHt}We^{-iHt},O]\|\le\varepsilon$ for all $|t|\le T$, then
    \begin{align}
        \|[\widetilde W,O]\| & \le\varepsilon
        +2\|W\|\,\|O\|\,e^{-(T\Delta E)^2/(2q)}.
        \label{eq:hastings_filtered_locality}
    \end{align}
    For $q=0$, $\widetilde W=W$ and $\|[W,O]\|\le\varepsilon$ from $t=0$
    \cite[lines 816--828]{Hastings2007AreaLaw}.
\end{theorem}
\begin{proof}\leanok
    Split (\ref{eq:hastings_gaussian_filtered_term}) at $|t|=T$. On $[-T,T]$
    the integrand norm is at most $\varepsilon$; on $\mathbb R\setminus[-T,T]$
    it is at most $2\|W\|\,\|O\|$, and with $h=q/\Delta E^2$ the substitution
    $t=T+s$ and $(T+s)^2\ge T^2+s^2$ bound the tail mass by $e^{-T^2/(2h)}$.
\end{proof}

\begin{theorem}[Localization of filtered operators by partial trace]
    \label{thm:hastings_filtered_partial_trace}
    \lean{
        TNLean.Lattice.norm_sub_siteExpectation_gaussianFilteredTerm_le_mass,
        TNLean.Lattice.norm_sub_siteExpectation_gaussianFilteredTerm_le
    }
    \leanok
    \uses{thm:hastings_filtered_locality}
    Let $\mathbb E_K$ be the normalized partial trace over $K^c$ tensored with
    $1_{K^c}$. If $q>0$, $T,\varepsilon\ge0$, and
    $\|e^{iHt}We^{-iHt}-\mathbb E_K(e^{iHt}We^{-iHt})\|\le\varepsilon$ for all
    $|t|\le T$, then $\mathbb E_K(\widetilde W)$ is supported on $K$, satisfies
    $\|\mathbb E_K(\widetilde W)\|\le\|W\|$, and
    \begin{align}
        \|\widetilde W-\mathbb E_K(\widetilde W)\| & \le\varepsilon
        +2\|W\|\,e^{-(T\Delta E)^2/(2q)}.
        \label{eq:hastings_filtered_partial_trace}
    \end{align}
    For $q=0$, $\|\widetilde W-\mathbb E_K(\widetilde W)\|\le\varepsilon$. This
    replaces the truncated-evolution approximation of
    \cite[lines 816--828]{Hastings2007AreaLaw}.
\end{theorem}
\begin{proof}\leanok
    Since $\mathrm{id}-\mathbb E_K$ is continuous and linear with
    $\|(\mathrm{id}-\mathbb E_K)(M)\|\le2\|M\|$, splitting the time integral at
    $|t|=T$ as in Theorem~\ref{thm:hastings_filtered_locality} gives
    (\ref{eq:hastings_filtered_partial_trace}).
\end{proof}

\subsection*{The localized parts and the cut-offs}

% The source prints the bound J^2 (line 794), which fails already for
% nearest-neighbour terms; the entry states the corrected constant.
\begin{theorem}[Commutators of the parts with the Hamiltonian]
    \label{thm:hastings_part_commutators}
    \lean{
        TNLean.MPS.AreaLaw.ChainHamiltonian.norm_commutator_leftPart_le,
        TNLean.MPS.AreaLaw.ChainHamiltonian.norm_commutator_rightPart_le,
        TNLean.MPS.AreaLaw.ChainHamiltonian.commutator_bufferPart_eq,
        TNLean.MPS.AreaLaw.ChainHamiltonian.norm_commutator_bufferPart_le,
        TNLean.MPS.AreaLaw.ChainHamiltonian.isSupportedOn_commutator_leftPart,
        TNLean.MPS.AreaLaw.ChainHamiltonian.isSupportedOn_commutator_rightPart
    }
    \leanok
    \uses{def:hastings_three_way_splitting}
    Let $H$ have range $R$ and term norm bound $J$ on $L$ sites, with
    three-way splitting $H_L,H_B,H_R$ at cut $j$ and buffer width $m$. With
    $c_R=2N_R(R+1)2^{2R+1}$, where $N_R$ bounds the number of terms crossing a
    cut,
    \begin{align}
        \|[H,H_L]\|,\ \|[H,H_R]\| & \le c_RJ^2,          &
        [H,H_B]                   & =-([H,H_L]+[H,H_R]), &
        \|[H,H_B]\|               & \le2c_RJ^2.
        \label{eq:hastings_part_commutators}
    \end{align}
    Moreover, $[H,H_L]$ is supported on
    $S_L=\{j-m-2R,\dots,j-m+R-1\}\cap\{0,\dots,L-1\}$ and $[H,H_R]$ on
    $S_R=\{j+m-R,\dots,j+m+2R-1\}\cap\{0,\dots,L-1\}$. The factor $c_R$
    corrects the printed bound $\|[H,H_L]\|\le J^2$ of
    \cite[line 794]{Hastings2007AreaLaw}; see
    \cite{gap:hastings07_commutator_constant}.
\end{theorem}

\begin{theorem}[Localized filtered parts]
    \label{thm:hastings_localized_parts}
    \lean{TNLean.MPS.AreaLaw.ChainHamiltonian.norm_gaussianFilteredTerm_leftPart_sub_localized_le,
        TNLean.MPS.AreaLaw.ChainHamiltonian.norm_gaussianFilteredTerm_rightPart_sub_localized_le,
        TNLean.MPS.AreaLaw.ChainHamiltonian.norm_gaussianFilteredTerm_bufferPart_sub_localized_le,
        TNLean.MPS.AreaLaw.ChainHamiltonian.isHermitian_leftPart_localized,
        TNLean.MPS.AreaLaw.ChainHamiltonian.isHermitian_rightPart_localized,
        TNLean.MPS.AreaLaw.ChainHamiltonian.isHermitian_bufferPart_localized,
        TNLean.MPS.AreaLaw.isSupportedOn_localized_sub_self}
    \leanok
    \uses{thm:hastings_filtered_partial_trace, thm:hastings_part_commutators,
        thm:chain_lieb_robinson, thm:hastings_gaussian_filtered_spectral,
        thm:hastings_three_way_splitting_sum}
    In the setting of Theorem~\ref{thm:hastings_part_commutators}, let $q>0$,
    $h=q/\Delta E^2$, $S_B=S_L\cup S_R$, and for $\ell\ge0$ let $K_\bullet$ be
    the $\ell$-neighbourhood of $S_\bullet$ ($\bullet\in\{L,B,R\}$). Define
    \begin{align}
        M_\bullet & =H_\bullet+\mathbb E_{K_\bullet}(\widetilde H_\bullet-H_\bullet).
        \label{eq:hastings_localized_part_def}
    \end{align}
    Then $M_\bullet=M_\bullet^\dagger$ and $M_\bullet-H_\bullet$ is supported
    on $K_\bullet$. With $(C',v)$ from Lemma~\ref{thm:chain_lieb_robinson} at
    $b_0=2^{2R+1}J$ and $c_R$ from Theorem~\ref{thm:hastings_part_commutators},
    for all $T\ge0$,
    \begin{align}
        \|\widetilde H_L-M_L\|,\ \|\widetilde H_R-M_R\|
         & \le T\,C'(3R)^2c_RJ^2e^{vT-\ell/2}
        +2c_RJ^2\sqrt{2h/\pi}\,e^{-T^2/(2h)},
        \label{eq:hastings_localized_parts_lr}   \\
        \|\widetilde H_B-M_B\|
         & \le T\,C'(6R)^2\,2c_RJ^2e^{vT-\ell/2}
        +4c_RJ^2\sqrt{2h/\pi}\,e^{-T^2/(2h)},
        \label{eq:hastings_localized_parts_b}
    \end{align}
    independently of $L$. This is
    \cite[lines 816--828]{Hastings2007AreaLaw} measured from
    $\widetilde H_\bullet$ rather than $H_\bullet$
    \cite{gap:hastings07_filtered_part_localization}. Since
    $\widetilde H_L+\widetilde H_B+\widetilde H_R=H-E_0$, the sum
    $M=M_L+M_B+M_R$ satisfies
    $\|M-(H-E_0)\|\le\sum_{\bullet\in\{L,B,R\}}\|\widetilde H_\bullet-M_\bullet\|$
    \cite[lines 766--767, 849--850]{Hastings2007AreaLaw}.
\end{theorem}

\begin{theorem}[The cut-offs $O_L$ and $O_R$]
    \label{thm:hastings_cutoffs}
    \uses{thm:hastings_localized_parts, thm:hastings_filtered_no_ground_component,
        def:hastings_spectral_cutoff, thm:hastings_spectral_cutoff_obd,
        thm:hastings_spectral_cutoff_norm, thm:hastings_spectral_cutoff_support,
        thm:hastings_spectral_cutoff_commute}
    For $\theta>0$ and Hermitian $M_L,M_R$ supported on $X_{1,j}$ and
    $X_{j+1,N}$, let $O_L=1_{[-\theta,\theta]}(M_L)$ and
    $O_R=1_{[-\theta,\theta]}(M_R)$. Then $O_L,O_R$ are commuting orthogonal
    projections supported on $X_{1,j}$ and $X_{j+1,N}$, and
    \begin{align}
        \|(O_L-1)\Psi_0\|       & \le\frac{\|M_L\Psi_0\|}{\theta}, &
        \|(O_R-1)\Psi_0\|       & \le\frac{\|M_R\Psi_0\|}{\theta}, &
        \|M_LO_L\|,\ \|M_RO_R\| & \le\theta.
        \label{eq:hastings_cutoffs_bounds}
    \end{align}
    For $q=(l/3)\Delta E/v$, buffer width and localization radius $\propto l$,
    and $\theta=\Delta E^{-1}J^2e^{-l/(6\xi)}$,
    Theorems~\ref{thm:hastings_filtered_no_ground_component}
    and~\ref{thm:hastings_localized_parts} yield
    $\|M_\bullet\Psi_0\|\le\mathcal O(\mathrm{poly}(l)e^{-l/(3\xi)})$ and
    $\|(O_\bullet-1)\Psi_0\|\le\mathcal O(\mathrm{poly}(l)e^{-l/(6\xi)})$
    \cite[lines 829--841]{Hastings2007AreaLaw}. Cutting off $|\lambda|\le\theta$
    rather than $\lambda\le\theta$ is required for
    $\|e^{i(M_L+M_R)t}O_LO_R-O_LO_R\|\le2|t|\theta$ at
    \cite[lines 879--880]{Hastings2007AreaLaw}; see
    \cite{gap:hastings07_spectral_cutoff_threshold}.
\end{theorem}
\begin{proof}
    Spectral projections of Hermitian operators on disjoint supports are
    commuting orthogonal projections on those supports;
    Theorems~\ref{thm:hastings_spectral_cutoff_obd}
    and~\ref{thm:hastings_spectral_cutoff_norm} give
    (\ref{eq:hastings_cutoffs_bounds}).
\end{proof}

\begin{theorem}[The cut-offs on the chain]
    \label{thm:hastings_cutoff_pair}
    \lean{TNLean.MPS.AreaLaw.ChainHamiltonian.exists_cutoffPair,
        TNLean.MPS.AreaLaw.ChainHamiltonian.localizedPart,
        TNLean.MPS.AreaLaw.ChainHamiltonian.isSupportedOn_localizedPart_leftPart,
        TNLean.MPS.AreaLaw.ChainHamiltonian.isSupportedOn_localizedPart_rightPart,
        TNLean.MPS.AreaLaw.ChainHamiltonian.norm_localizedPart_apply_le,
        TNLean.MPS.AreaLaw.partLocalizationError,
        TNLean.MPS.AreaLaw.leftWindow, TNLean.MPS.AreaLaw.rightWindow,
        TNLean.MPS.AreaLaw.chainNeighborhood_interval_subset,
        TNLean.Lattice.norm_gaussianFilteredTerm_sub_self_le}
    \leanok
    \uses{thm:hastings_localized_parts, thm:hastings_cutoffs,
        thm:hastings_filtered_no_ground_component, thm:hastings_part_commutators}
    Let $\Psi_0$ be a gapped ground vector of $H$ (range $R$, term norm bound
    $J\ge0$, ground energy $E_0$, gap $\Delta E>0$). For cut $j$, buffer width
    $m$, filter parameter $q\ge0$, and localization radius $k$ with
    $R+k\le m$, let $M_L,M_R$ be the localized centered parts of
    Theorem~\ref{thm:hastings_localized_parts} ($\ell=k$), $O_L,O_R$ their
    cut-offs at $\theta>0$, and $H_0=M_L+M_R$. Then $M_L,O_L$ act on $X_{1,j}$,
    $M_R,O_R$ act on $X_{j+1,N}$, and for every $T\ge0$,
    \begin{align}
        \|\Psi_0-O_LO_R\Psi_0\|
         & \le\frac2\theta\bigl(\Delta E^{-1}e^{-q/2}c_RJ^2+\delta\bigr), &
        \|H_0O_LO_R\|
         & \le2\theta,
        \label{eq:hastings_cutoff_pair}
    \end{align}
    where $\delta$ is the right-hand side of
    (\ref{eq:hastings_localized_parts_lr}) at $\ell=k$. Moreover,
    $H-E_0-H_0$ is supported within distance $m+2R+k$ of $j$ and satisfies,
    uniformly in $L$,
    \begin{align}
        \|H-E_0-H_0\| & \le4(m+R)2^{2R+1}J+2\sqrt{2h/\pi}\,c_RJ^2,
        \qquad h=q/\Delta E^2
        \label{eq:hastings_cutoff_pair_remainder}
    \end{align}
    \cite[lines 766--842]{Hastings2007AreaLaw}.
\end{theorem}
\begin{proof}\leanok
    By Theorems~\ref{thm:hastings_filtered_no_ground_component}
    and~\ref{thm:hastings_localized_parts},
    $\|M_\bullet\Psi_0\|\le\|\widetilde H_\bullet\Psi_0\|
        +\|\widetilde H_\bullet-M_\bullet\|\le\Delta E^{-1}e^{-q/2}c_RJ^2+\delta$,
    so Theorem~\ref{thm:hastings_cutoffs} gives (\ref{eq:hastings_cutoff_pair}).
    Since $H-E_0-H_0=H_B-\sum_{\bullet\in\{L,R\}}\mathbb E_{K_\bullet}(\widetilde H_\bullet-H_\bullet)$,
    using $\|H_B\|\le4(m+R)2^{2R+1}J$ and the Duhamel bound
    $\|\widetilde H_\bullet-H_\bullet\|\le\sqrt{2h/\pi}\,\|[H,H_\bullet]\|$
    yields (\ref{eq:hastings_cutoff_pair_remainder}).
\end{proof}

\subsection*{The Gaussian time average and the boundary operator}

For $q>0$, the Gaussian time average $P_q$ of
Definition~\ref{def:hastings_gaussian_time_average} with $E_0=0$ and
$M=M_L+M_B+M_R=M^\dagger$ satisfies $\|P_q(H)-P_0\|\le e^{-q/2}$ and
$\|P_q(M)-P_q(H)\|\le\sqrt q\,\Delta E^{-1}\|M-H\|$ by
Theorems~\ref{thm:hastings_gaussian_time_average_gap}
and~\ref{thm:hastings_gaussian_time_average_stability}
\cite[lines 843--864]{Hastings2007AreaLaw}. If $\|M-H\|\le\delta$ and
$\|\Psi_0-O_LO_R\Psi_0\|\le\eta$, Theorem~\ref{thm:hastings_moved_estimate}
gives $\|P_q(M)O_LO_R-P_0\|\le e^{-q/2}+\sqrt q\,\Delta E^{-1}\delta+\eta$
\cite[lines 865--876]{Hastings2007AreaLaw}. Writing
$e^{i(A+B)t}=U(t)e^{iAt}$ and averaging $U(t)$ to $P_B$
(Definitions~\ref{def:hastings_interaction_propagator}
and~\ref{def:hastings_interaction_average}) with $A=M_L+M_R$, $B=M_B$, and
$X=O_LO_R$, Theorem~\ref{thm:hastings_interaction_average_last} bounds
$\|P_BX-P_0\|$ \cite[lines 877--905]{Hastings2007AreaLaw}.

\begin{theorem}[Localization of the boundary average]
    \label{thm:hastings_boundary_localization}
    \lean{TNLean.MPS.AreaLaw.Hastings.exists_supported_interactionAverage_approx,
        TNLean.MPS.AreaLaw.Hastings.exists_supported_interactionAverage_approx_exp}
    \leanok
    \uses{def:hastings_interaction_average, thm:chain_lieb_robinson}
    Let $H$ have range $R$, local dimension $d\ge1$, and term norm bound $J$ on
    $L$ sites. Let $H_0,A,B$ be Hermitian with $H-H_0$ supported on $Z$, and let
    $P_B$ be the interaction average of
    Definition~\ref{def:hastings_interaction_average} for $(A,B,h)$ with $h>0$.
    With $(C',v)$ from Lemma~\ref{thm:chain_lieb_robinson} at
    $b_0=2^{2R+1}J$, for every $\ell,T\ge0$ there exists $O_B$ supported within
    distance $\ell$ of $Z$ such that $\|O_B\|\le1$ and
    \begin{align}
        \|P_B-O_B\| & \le\sqrt h\,\bigl(\|A+B-H\|+\|A-H_0\|
        +C'|Z|^2\|H-H_0\|e^{vT-\ell/2}\bigr)+4e^{-T^2/(2h)}.
        \label{eq:hastings_boundary_localization}
    \end{align}
    At $T=\ell/(4v)$, the last two terms become
    $\sqrt h\,C'|Z|^2\|H-H_0\|e^{-\ell/4}+4e^{-\ell^2/(32v^2h)}$.
    No positivity of $O_B$ is claimed \cite[note added]{Hastings2007AreaLaw}.
\end{theorem}
\begin{proof}\leanok
    \uses{thm:qc_supported_projected_relative_evolution}
    The relative evolution $V(t)=e^{iHt}e^{-iH_0t}$ has generator
    $e^{iHs}(H-H_0)e^{-iHs}$, which Lemma~\ref{thm:chain_lieb_robinson}
    approximates within distance $\ell$ of $Z$ up to
    $C'|Z|^2\|H-H_0\|e^{v|s|-\ell/2}$. By
    Theorem~\ref{thm:qc_supported_projected_relative_evolution}, the projected
    generator integrates to a unitary $\widetilde V(t)$ supported there with
    $\|V(t)-\widetilde V(t)\|\le|t|C'|Z|^2\|H-H_0\|e^{vT-\ell/2}$ for
    $|t|\le T$, while $\|U(t)-V(t)\|\le|t|(\|A+B-H\|+\|A-H_0\|)$. Setting
    $O_B=\int\widetilde V(t)\,d\gamma_h(t)$ and using
    $\int|t|\,d\gamma_h\le\sqrt h$ on $[-T,T]$ and $\|U(t)-\widetilde V(t)\|\le2$
    with tail mass $2e^{-T^2/(2h)}$ on $\mathbb R\setminus[-T,T]$ gives
    (\ref{eq:hastings_boundary_localization}).
\end{proof}

With $A=H_0+E_0$ and $B=H-E_0-H_0$ (which leaves $P_B$ unchanged since
$e^{\pm iE_0t}$ cancel), $\|A-H_0\|=\|A+B-H\|=0$, $|Z|=\mathcal O(l)$, and
$h=q/\Delta E^2$ with $q,\ell\propto l$, so
(\ref{eq:hastings_boundary_localization}) yields the
$\mathcal O(\Delta E^{-1}J\sqrt q\,e^{-l/(6\xi)})$ bound of
\cite[lines 896--905]{Hastings2007AreaLaw} with $O_B$ supported on
$X_{j-l+1,j+l}$.

\subsection*{The factorization and its positive version}

\begin{theorem}[Approximate factorization without positivity]
    \label{thm:hastings_projector_factorization}
    \lean{TNLean.MPS.AreaLaw.exists_projectorFactorization,
        TNLean.MPS.AreaLaw.exists_projectorFactorization_posSemidef,
        TNLean.MPS.AreaLaw.exists_isFactorizationAt_sqrt_decay,
        TNLean.MPS.AreaLaw.lemmaOneLength,
        TNLean.MPS.AreaLaw.Hastings.lemmaOne_errorBudget,
        TNLean.MPS.AreaLaw.Hastings.lemmaOneBudgetA,
        TNLean.MPS.AreaLaw.Hastings.lemmaOneBudgetB,
        TNLean.MPS.AreaLaw.lemmaOneSqrtA,
        TNLean.MPS.AreaLaw.lemmaOneSqrtB,
        TNLean.MPS.AreaLaw.lemmaOneConst,
        TNLean.MPS.AreaLaw.lemmaOneConst_nonneg,
        TNLean.MPS.AreaLaw.exists_isFactorizationAt_lemmaOneConst,
        TNLean.MPS.AreaLaw.isFactorizationAt_zero_one_one,
        Matrix.IsGappedGroundState.sub_smul_one}
    \leanok
    \uses{thm:hastings_three_way_splitting_sum,
        thm:hastings_three_way_splitting_supports,
        thm:hastings_filtered_no_ground_component,
        thm:hastings_part_commutators, thm:hastings_localized_parts,
        thm:hastings_cutoffs, thm:hastings_gaussian_time_average_stability,
        thm:hastings_moved_estimate, thm:hastings_interaction_average_last,
        thm:hastings_boundary_localization}
    Let $H$ have range $R$, term norm bound $J$, unique ground state $\Psi_0$,
    gap $\Delta E>0$, and ground projector
    $P_0=|\Psi_0\rangle\langle\Psi_0|$. Set $\xi=8\max(2v_0/\Delta E,1)$, where
    $v_0$ is the chain localization velocity, and choose
    \begin{align}
        q & =\frac{2l}{3\xi},
        \label{eq:hastings_lemma1_q}
    \end{align}
    which doubles the printed value of \cite[line 792]{Hastings2007AreaLaw} so
    that $e^{-q/2}=e^{-l/(3\xi)}$ as used in
    \cite[line 848]{Hastings2007AreaLaw}
    \cite{gap:hastings07_gaussian_average_exponent}. For every $\xi'>6\xi$
    there exists $C_1=C_1(R,J,\Delta E,\xi'-6\xi)\ge0$, independent of $L$, the
    local dimension, and $\Psi_0$, such that for every cut $j$ and every $l$
    there exist $O_B,O_L,O_R$ satisfying:
    \begin{enumerate}
        \item $O_L,O_R$ are commuting orthogonal projections supported on
              $X_{1,j}$ and $X_{j+1,N}$;
        \item $\|O_B\|\le1$ and $O_B$ is supported on $X_{j-l+1,j+l}$;
        \item $\|O_BO_LO_R-P_0\|\le C_1e^{-l/\xi'}$.
    \end{enumerate}
    Every $\xi'>6\xi$ absorbs the $\sqrt q\propto\sqrt l$ prefactor of
    \cite[lines 903--907]{Hastings2007AreaLaw}, whereas $\xi'=6\xi$
    \cite[lines 124--144]{Hastings2007AreaLaw} and positivity of $O_B$
    \cite[lines 912--913]{Hastings2007AreaLaw} are not claimed, and $(v_0,2)$
    is not shown to be a Lieb--Robinson pair in the source's sense
    \cite{gap:hastings07_lemma1_decay_length}. Replacing $O_B$ by
    $O_B^\dagger O_B\ge0$ gives rate $e^{-l/(2\xi')}$ by
    Theorem~\ref{thm:hastings_positive_factorization}.
\end{theorem}
\begin{proof}\leanok
    \uses{thm:hastings_cutoff_pair, thm:hastings_moved_estimate,
        thm:hastings_boundary_factorization, thm:hastings_decay_slack,
        thm:hastings_positive_factorization, thm:hastings_factorization_rate,
        thm:hastings_positive_factorization_rate}
    Shift $H$ so that $E_0=0$. For $l\ge12R+8$, set $s=2l/(3\xi)$,
    $\theta=e^{-s/4}$, $T=s/\Delta E$, $k=\lceil l/4\rceil$, $m=R+k$, and
    $\ell'=l-3R-2k$. Combining Theorems~\ref{thm:hastings_cutoff_pair},
    \ref{thm:hastings_moved_estimate}, and
    \ref{thm:hastings_boundary_factorization} produces $O_L,O_R,O_B$ with $O_B$
    supported within $m+2R+k+\ell'=l$ of $j$ and total error bounded by
    $(a+b\sqrt l)e^{-l/(6\xi)}$ for $a,b$ depending only on $(R,J,\Delta E)$.
    For $l<12R+8$, take $O_B=0$ and $O_L=O_R=1$. Theorem~\ref{thm:hastings_decay_slack}
    absorbs $\sqrt l$ into $C_1e^{-l/\xi'}$ for $\xi'>6\xi$.
\end{proof}

\begin{theorem}[Approximate factorization in units of the gap]
    \label{thm:hastings_projector_factorization_scaled}
    \lean{TNLean.MPS.AreaLaw.exists_projectorFactorization_div,
        TNLean.MPS.AreaLaw.lemmaOneLength_div,
        TNLean.MPS.AreaLaw.ChainHamiltonian.smul,
        Matrix.IsGappedGroundState.smul,
        Entropy.IsSupportedOn.smul}
    \leanok
    \uses{thm:hastings_projector_factorization}
    Under the assumptions of Theorem~\ref{thm:hastings_projector_factorization},
    set $c_R=2(R+1)2^{2R+1}e^{2R}$, $v=8c_RJ$, and
    \begin{align}
        \hat\xi & =16\Bigl(1+\frac{c_RJ}{\Delta E}\Bigr)
        \le2\max\Bigl(\frac{2v}{\Delta E},16\Bigr).
        \label{eq:hastings_lemma1_hat_xi}
    \end{align}
    For every $\xi'>6\hat\xi$ there exists $C_1=C_1(R,J/\Delta E,\xi')\ge0$
    such that for every $j$ and $l$ there exist $O_B,O_L,O_R$ satisfying items
    1--2 of Theorem~\ref{thm:hastings_projector_factorization} and
    $\|O_BO_LO_R-P_0\|\le C_1e^{-l/\xi'}$
    \cite[lines 96--112, 139--142]{Hastings2007AreaLaw},
    \cite{gap:hastings07_lemma1_decay_length}.
\end{theorem}
\begin{proof}\leanok
    \uses{thm:hastings_projector_factorization}
    Apply Theorem~\ref{thm:hastings_projector_factorization} to $H/\Delta E$,
    which has range $R$, term norm bound $J/\Delta E$, gap $1$, and length
    $8\max(2(1+c_RJ/\Delta E),1)=\hat\xi$.
\end{proof}

\begin{theorem}[A polynomial bound on the constant of the factorization]
    \label{thm:hastings_lemma1_polynomial_constant}
    \lean{TNLean.MPS.AreaLaw.exists_projectorFactorization_pow,
        TNLean.MPS.AreaLaw.exists_projectorFactorization_posSemidef_pow,
        TNLean.MPS.AreaLaw.lemmaOneConst_le_pow,
        TNLean.MPS.AreaLaw.sixteen_le_lemmaOneLength_and_le}
    \leanok
    \uses{thm:hastings_projector_factorization_scaled,
        thm:hastings_positive_factorization}
    For every $\kappa>6$ there exists $K=K(R,\kappa)\ge0$ such that, under the
    assumptions of Theorem~\ref{thm:hastings_projector_factorization} with
    $\hat\xi$ as in (\ref{eq:hastings_lemma1_hat_xi}), for every $j$ and $l$
    there exist $O_B,O_L,O_R$ satisfying items 1--2 of
    Theorem~\ref{thm:hastings_projector_factorization} and
    \begin{align}
        \|O_BO_LO_R-P_0\| & \le K\hat\xi^3e^{-l/(\kappa\hat\xi)}.
        \label{eq:hastings_lemma1_polynomial_constant}
    \end{align}
    Replacing $O_B$ by $O_B^\dagger O_B\ge0$ gives rate
    $e^{-l/(2\kappa\hat\xi)}$ with another constant $K(R,\kappa)$
    \cite[lines 258--262]{Hastings2007AreaLaw},
    \cite{gap:hastings07_lemma1_positivity_rate,
        gap:hastings07_lemma1_decay_length,
        gap:hastings07_gaussian_average_exponent}.
\end{theorem}
\begin{proof}\leanok
    \uses{thm:hastings_projector_factorization_scaled,
        thm:hastings_projector_factorization,
        thm:hastings_positive_factorization_rate,
        thm:hastings_factorization_mono}
    In Theorem~\ref{thm:hastings_projector_factorization} applied to
    $H/\Delta E$ at $\xi'=\kappa\hat\xi$ with $x=J/\Delta E\le\hat\xi$ and
    $\hat\xi\ge16$, the budget parameters satisfy $a=\mathcal O_R(\hat\xi^3)$,
    $b=\mathcal O_R(\hat\xi^2)$, and
    $\bigl(2e(1/(6\hat\xi)-1/(\kappa\hat\xi))\bigr)^{-1/2}
        =\sqrt{\hat\xi/\beta}\le\hat\xi/\sqrt\beta$ for
    $\beta=2e(\kappa-6)/(6\kappa)$, giving
    $C_1=a+b\sqrt{\hat\xi/\beta}\le K\hat\xi^3$.
\end{proof}

\begin{theorem}[The factorization at the printed rate, for a raised velocity]
    \label{thm:hastings_lemma1_endpoint}
    \lean{TNLean.MPS.AreaLaw.exists_projectorFactorization_endpoint,
        TNLean.MPS.AreaLaw.exists_projectorFactorization_posSemidef_endpoint,
        TNLean.MPS.AreaLaw.lemmaOneEndpointLength,
        TNLean.MPS.AreaLaw.lemmaOneEndpointLength_eq}
    \leanok
    \uses{thm:hastings_lemma1_polynomial_constant}
    With $v_0=1+c_RJ/\Delta E$, $\hat\xi=8\max(2v_0,1)$, and raised velocity
    $v'=\tfrac54v_0$, set $\tilde\xi=8\max(2v',1)=\tfrac54\hat\xi$. There
    exists $K=K(R)\ge0$ such that, under the assumptions of
    Theorem~\ref{thm:hastings_projector_factorization}, for every $j$ and $l$
    there exist $O_B,O_L,O_R$ satisfying items 1--2 of
    Theorem~\ref{thm:hastings_projector_factorization} and
    \begin{align}
        \|O_BO_LO_R-P_0\| & \le K\tilde\xi^3e^{-l/(6\tilde\xi)},
        \label{eq:hastings_lemma1_endpoint}
    \end{align}
    and with $O_B\ge0$ at rate $e^{-l/(12\tilde\xi)}$
    \cite[lines 96--112, 141--142]{Hastings2007AreaLaw},
    \cite{gap:hastings07_lemma1_decay_length}.
\end{theorem}
\begin{proof}\leanok
    \uses{thm:hastings_lemma1_polynomial_constant}
    Apply Theorem~\ref{thm:hastings_lemma1_polynomial_constant} at
    $\kappa=15/2$, using $6\tilde\xi=\tfrac{15}2\hat\xi$ and
    $\hat\xi^3\le\tilde\xi^3$.
\end{proof}

\begin{theorem}[An effect operator from an approximate factorization]
    \label{thm:hastings_effect_from_factorization}
    \lean{Matrix.norm_mul_conjTranspose_mul_self_mul_sub_le}
    \leanok
    For any matrix $B$ and orthogonal projections $P,Q$ with
    $\|BQ-P\|\le\varepsilon$,
    \begin{align}
        \|QB^\dagger BQ-P\| & \le2\varepsilon+\varepsilon^2
        \label{eq:hastings_effect_from_factorization}
    \end{align}
    \cite[lines 916--925]{Hastings2007AreaLaw}.
\end{theorem}
\begin{proof}\leanok
    Put $E=BQ-P$. Since $Q=Q^\dagger$ and $P^2=P=P^\dagger$,
    $QB^\dagger BQ-P=(P+E^\dagger)(P+E)-P=E^\dagger P+PE+E^\dagger E$, and
    $\|P\|\le1$ gives (\ref{eq:hastings_effect_from_factorization}).
\end{proof}

By Theorem~\ref{thm:hastings_effect_from_factorization} with $B=O_B$, $P=P_0$,
and $Q=O_LO_R$, the positive operator $O_B^\dagger O_B$ suffices for all
ground-state and marginal expectation estimates in the area-law proof at rate
$e^{-l/\xi'}$ ($\xi'>6\xi$), whereas the direct factorization with
$O_B^+=O_B^\dagger O_B\ge0$ has error $\mathcal O(\sqrt{\varepsilon(l)})$ and
rate $e^{-l/(2\xi')}$ by Theorem~\ref{thm:hastings_positive_repair}
\cite[lines 932--940]{Hastings2007AreaLaw}.

% Lean statement pending.
\begin{lemma}[Approximate factorization of the ground-state projector]
    \label{lem:hastings_lemma1}
    \uses{thm:hastings_projector_factorization,
        thm:hastings_effect_from_factorization,
        thm:hastings_positive_factorization,
        thm:hastings_lemma1_polynomial_constant,
        thm:hastings_lemma1_endpoint}
    Let $H$ satisfy the finite-range, finite-strength, and gap conditions of
    \cite[Lemma~1, lines 246--267]{Hastings2007AreaLaw}. For any $j$ and $l$
    there exist Hermitian positive definite $O_B(j,l),O_L(j,l),O_R(j,l)$ of
    norm at most $1$, supported on $X_{j-l+1,j+l}$, $X_{1,j}$, and
    $X_{j+1,N}$, such that for $P_0=|\Psi_0\rangle\langle\Psi_0|$ and a
    polynomial $C_1(\xi)$,
    \begin{align}
        \|O_B(j,l)O_L(j,l)O_R(j,l)-P_0\| & \le C_1(\xi)e^{-l/\xi'}
        \equiv\varepsilon(l),\qquad\xi'=6\xi.
        \notag
    \end{align}
\end{lemma}

As recorded in Theorems~\ref{thm:hastings_projector_factorization}--\ref{thm:hastings_lemma1_endpoint}
and Theorem~\ref{thm:hastings_positive_repair}, the appendix proof yields the
non-positive factorization at $\xi'>6\xi$ (or $\xi'=6\tilde\xi$ for the raised
length $\tilde\xi=\tfrac54\hat\xi$) and the positive semidefinite
replacement $O_B^\dagger O_B$ at rate $e^{-l/(2\xi')}$, rather than
Lemma~\ref{lem:hastings_lemma1} as printed.
